% GENERATED FILE. Edit canonical lessons, then run npm run generate:books.
% canonical-source-hash: fnv1a64:7b9a5613f08b052b
% canonical-lessons: TA-C200-killu, TA-C200-valai, TA-C200-pucu, TA-C200-cekari, TA-C200-urincu

\chapter{To pinch, To bend, To paint, To collect, To suck}
\label{ch:ta-to-pinch-to-bend-to-paint-to-collect-to-suck}
\hlchaptermodality{200}

\begin{chapteropening}
\textbf{By the end of this chapter:} \emph{I can say to pinch, to bend, to paint, to collect and to suck.}

Complete the last lesson of chapter 200: I can say to pinch, to bend, to paint, to collect and to suck.
\end{chapteropening}

\section[\texorpdfstring{\ta{கிள்ளு} (kiḷḷu)}{கிள்ளு (kiḷḷu)}]{\ta{கிள்ளு} (kiḷḷu) — to pinch, to nip off}

\label{lesson:TA-C200-killu}



\begin{quote}
Before the new one: say the Tamil for to discuss, then the Tamil for to compare.
\end{quote}

\subsection*{You'll want to know: \ta{கிள்ளு}}
\textbf{\ta{கிள்ளு}} — \emph{kiḷḷu} — \textquotedblleft{}to pinch, to nip off\textquotedblright{}.

\begin{grammarlens}[title={What you've built}]
One more word for this chapter.
\end{grammarlens}

\subsection*{Your turn}
\begin{itemize}
\raggedright
  \item \emph{Say it:} \emph{kiḷḷu}
  \item \emph{Say it:} \emph{kiḷḷu}, once more
  \item \emph{Say it:} say \emph{kiḷḷu}
  \item \emph{Recall:} say the Tamil for to discuss, then the Tamil for to compare, then say \emph{kiḷḷu} again
\end{itemize}

\begin{tcolorbox}[breakable,colback=teal!4,colframe=teal!35!black,title={Before you move on}]
What does \ta{கிள்ளு} mean? (\textquotedblleft{}To pinch, to nip off\textquotedblright{}.) Say it once more.
\end{tcolorbox}
\section[\texorpdfstring{\ta{வளை} (vaḷai)}{வளை (vaḷai)}]{\ta{வளை} (vaḷai) — to bend}

\label{lesson:TA-C200-valai}



\begin{quote}
Before the new one: say the Tamil for to compare, then the Tamil for to pinch.
\end{quote}

\subsection*{You'll want to know: \ta{வளை}}
\textbf{\ta{வளை}} — \emph{vaḷai} — \textquotedblleft{}to bend\textquotedblright{}.

\begin{grammarlens}[title={What you've built}]
One more word for this chapter.
\end{grammarlens}

\subsection*{Your turn}
\begin{itemize}
\raggedright
  \item \emph{Say it:} \emph{vaḷai}
  \item \emph{Say it:} \emph{vaḷai}, once more
  \item \emph{Say it:} say \emph{vaḷai}
  \item \emph{Recall:} say the Tamil for to compare, then the Tamil for to pinch, then say \emph{vaḷai} again
\end{itemize}

\begin{tcolorbox}[breakable,colback=teal!4,colframe=teal!35!black,title={Before you move on}]
What does \ta{வளை} mean? (\textquotedblleft{}To bend\textquotedblright{}.) Say it once more.
\end{tcolorbox}
\section[\texorpdfstring{\ta{பூசு} (pūcu)}{பூசு (pūcu)}]{\ta{பூசு} (pūcu) — to paint (a wall), to coat, to smear on}

\label{lesson:TA-C200-pucu}



\begin{quote}
Before the new one: say the Tamil for to pinch, then the Tamil for to bend.
\end{quote}

\subsection*{You'll want to know: \ta{பூசு}}
\textbf{\ta{பூசு}} — \emph{pūcu} — \textquotedblleft{}to paint (a wall), to coat, to smear on\textquotedblright{}.

\begin{grammarlens}[title={What you've built}]
One more word for this chapter.
\end{grammarlens}

\subsection*{Your turn}
\begin{itemize}
\raggedright
  \item \emph{Say it:} \emph{pūcu}
  \item \emph{Say it:} \emph{pūcu}, once more
  \item \emph{Say it:} say \emph{pūcu}
  \item \emph{Recall:} say the Tamil for to pinch, then the Tamil for to bend, then say \emph{pūcu} again
\end{itemize}

\begin{tcolorbox}[breakable,colback=teal!4,colframe=teal!35!black,title={Before you move on}]
What does \ta{பூசு} mean? (\textquotedblleft{}To paint (a wall), to coat, to smear on\textquotedblright{}.) Say it once more.
\end{tcolorbox}
\section[\texorpdfstring{\ta{சேகரி} (cēkari)}{சேகரி (cēkari)}]{\ta{சேகரி} (cēkari) — to collect, to gather}

\label{lesson:TA-C200-cekari}



\begin{quote}
Before the new one: say the Tamil for to bend, then the Tamil for to paint.
\end{quote}

\subsection*{You'll want to know: \ta{சேகரி}}
\textbf{\ta{சேகரி}} — \emph{cēkari} — \textquotedblleft{}to collect, to gather\textquotedblright{}.

\begin{grammarlens}[title={What you've built}]
One more word for this chapter.
\end{grammarlens}

\subsection*{Your turn}
\begin{itemize}
\raggedright
  \item \emph{Say it:} \emph{cēkari}
  \item \emph{Say it:} \emph{cēkari}, once more
  \item \emph{Say it:} say \emph{cēkari}
  \item \emph{Recall:} say the Tamil for to bend, then the Tamil for to paint, then say \emph{cēkari} again
\end{itemize}

\begin{tcolorbox}[breakable,colback=teal!4,colframe=teal!35!black,title={Before you move on}]
What does \ta{சேகரி} mean? (\textquotedblleft{}To collect, to gather\textquotedblright{}.) Say it once more.
\end{tcolorbox}
\section[\texorpdfstring{\ta{உறிஞ்சு} (uṟiñcu)}{உறிஞ்சு (uṟiñcu)}]{\ta{உறிஞ்சு} (uṟiñcu) — to suck, to sip through a straw}

\label{lesson:TA-C200-urincu}



\begin{quote}
Before the new one: say the Tamil for to paint, then the Tamil for to collect.
\end{quote}

\subsection*{You'll want to know: \ta{உறிஞ்சு}}
\textbf{\ta{உறிஞ்சு}} — \emph{uṟiñcu} — \textquotedblleft{}to suck, to sip through a straw\textquotedblright{}.

\begin{grammarlens}[title={What you've built}]
One more word for this chapter.
\end{grammarlens}

\subsection*{Your turn}
\begin{itemize}
\raggedright
  \item \emph{Say it:} \emph{uṟiñcu}
  \item \emph{Say it:} \emph{uṟiñcu}, once more
  \item \emph{Say it:} say \emph{uṟiñcu}
  \item \emph{Recall:} say the Tamil for to paint, then the Tamil for to collect, then say \emph{uṟiñcu} again
\end{itemize}

\begin{tcolorbox}[breakable,colback=teal!4,colframe=teal!35!black,title={Before you move on}]
What does \ta{உறிஞ்சு} mean? (\textquotedblleft{}To suck, to sip through a straw\textquotedblright{}.) Say it once more.
\end{tcolorbox}
